\documentclass[aspectratio=169, xcolor={table,xcdraw}]{beamer}

\usepackage[T1]{fontenc}
\usepackage[utf8]{inputenc}
\usepackage{tgpagella}
\usepackage{tgcursor}
\usefonttheme{serif}
\usepackage[spanish, es-noshorthands]{babel}
\renewcommand{\tablename}{Tabla}

\usepackage{amsmath, amssymb, bm}
\usepackage{graphicx}
\usepackage{booktabs}
\usepackage{array} % Requerido para raggedright en tablas p{}
\usepackage[most]{tcolorbox}
\usepackage{tikz}
\usetikzlibrary{shapes.geometric, arrows.meta, positioning, calc}

\usetheme{Madrid}
\usecolortheme{default}
\definecolor{ucbblue}{RGB}{0, 51, 102}
\setbeamercolor{structure}{fg=ucbblue}
\setbeamercolor{title}{fg=white, bg=ucbblue}
\setbeamercolor{frametitle}{fg=white, bg=ucbblue}

\title[Integración Unidad II]{Sesión de Integración y Diagnóstico (Pre-EC2)}
\subtitle{De la Pose Deseada a la Dinámica del Actuador}
\author[B. Quiroga Turdera]{M.Sc. Ing. Bernardo Quiroga Turdera}
\institute[UCB Tarija]{Universidad Católica Boliviana ``San Pablo'' — Sede Tarija}
\date{Semana 10 - Sesión 1}

\begin{document}

\begin{frame}
    \titlepage
\end{frame}

\begin{frame}{Objetivo de la Sesión: Diagnóstico y Preparación EC2[cite: 15]}
    \begin{tcolorbox}[colback=white, colframe=ucbblue, title=Alineación Curricular{[cite: 15]}, boxsep=1mm]
        \small
        El objetivo de hoy no es introducir nuevas fórmulas, sino \textbf{evaluar y consolidar su capacidad de ejecución independiente} sobre las competencias clave de la Unidad II antes del EC2[cite: 15].
    \end{tcolorbox}
    
    \vfill
    \textbf{Pipeline de la Unidad II (El flujo completo)[cite: 15]:}
    \begin{center}
    \begin{tikzpicture}[
        node distance=0.4cm and 0.5cm,
        box/.style={rectangle, draw=ucbblue, thick, fill=ucbblue!5, text width=2.1cm, align=center, rounded corners, font=\scriptsize},
        arrow/.style={-Stealth, thick, ucbblue}
    ]
        \node[box] (ik) {\textbf{Pieper}\\$T_d \to \mathbf{q}_f$};
        \node[box, right=of ik] (lspb) {\textbf{LSPB}\\$q(t), \dot{q}(t), \ddot{q}(t)$};
        
        \node[box, right=of lspb] (jac) {\textbf{Jacobiano}\\$J(\mathbf{q})$ $\to$ $\dot{\mathbf{x}}(t)$};
        \node[box, below=0.4cm of jac] (sing) {\textbf{Singularidad}\\Rango y $\det$};
        \node[box, left=of sing] (dyn) {\textbf{Dinámica}\\$M, C, g \to \tau(t)$};
        
        \draw[arrow] (ik) -- (lspb);
        \draw[arrow] (lspb) -- (jac);
        \draw[arrow] (jac) -- (sing);
        \draw[arrow] (lspb) -- (dyn);
    \end{tikzpicture}
    \end{center}
    \vfill
    \footnotesize \textit{Temas de control de espacio nulo, redundancia y matriz pseudoinversa se mantienen como material diferido para las siguientes etapas[cite: 15].}
\end{frame}

\begin{frame}{Selección de Herramientas de Ingeniería[cite: 15]}
    \begin{center}
        \renewcommand{\arraystretch}{1.5}
        \footnotesize
        \begin{tabular}{|>{\raggedright\arraybackslash}p{7cm}|>{\raggedright\arraybackslash}p{6cm}|}
        \hline
        \rowcolor{ucbblue!20} \textbf{Pregunta de Ingeniería} & \textbf{Herramienta / Modelo Principal} \\ \hline
        ¿Qué configuraciones articulares logran una pose deseada? & Cinemática Inversa Analítica (Pieper)[cite: 15] \\ \hline
        ¿El candidato matemático realmente reproduce la pose? & Cinemática Directa (FK)[cite: 15] \\ \hline
        ¿Qué velocidad cartesiana resulta de la velocidad articular? & Jacobiano Geométrico[cite: 15] \\ \hline
        ¿Dónde pierde el robot movilidad instantánea? & Rango del Jacobiano (Singularidad)[cite: 15] \\ \hline
        ¿Cómo deben moverse las juntas continuamente en el tiempo? & Generación de Trayectorias (LSPB)[cite: 15] \\ \hline
        ¿Qué torque es requerido para lograr $q, \dot{q}, \ddot{q}$? & Dinámica (Euler-Lagrange)[cite: 15] \\ \hline
        \end{tabular}
    \end{center}
\end{frame}

\begin{frame}{Dinámica de Trabajo de Hoy[cite: 15]}
    Seguiremos el ciclo diagnóstico formal[cite: 15]:
    \begin{center}
    \begin{tikzpicture}[
        node distance=0.4cm,
        box/.style={rectangle, draw=red!80!black, thick, fill=red!5, text width=1.9cm, align=center, rounded corners, font=\scriptsize},
        arrow/.style={-Stealth, thick, red!80!black}
    ]
        \node[box] (att) {\textbf{1. Intento}\\Independiente};
        \node[box, right=of att] (diag) {\textbf{2. Diagnóstico}\\de Errores};
        \node[box, right=of diag] (corr) {\textbf{3. Corrección}\\Focalizada};
        \node[box, right=of corr] (app) {\textbf{4. Segunda}\\Aplicación};
        \node[box, fill=green!10, right=of app] (read) {\textbf{5. Juicio}\\de Readiness};
        
        \draw[arrow] (att) -- (diag);
        \draw[arrow] (diag) -- (corr);
        \draw[arrow] (corr) -- (app);
        \draw[arrow] (app) -- (read);
    \end{tikzpicture}
    \end{center}
    
    \vspace{0.4cm}
    \textbf{Tienen 5 tareas que resolver en su hoja de ejercicios[cite: 15]:}
    \begin{enumerate}
        \item Cinemática Inversa (Pieper).
        \item Jacobiano y Singularidades.
        \item Generación de Trayectoria LSPB.
        \item Dinámica Euler-Lagrange.
        \item Integración Sistémica.
    \end{enumerate}
    \textit{Trabajen el primer intento individualmente y sin apuntes para medir su nivel real.}
\end{frame}

\end{document}
